\documentclass[12pt]{book} % font 10
%\usepackage[T1]{fontenc}
\usepackage[a4paper]{geometry}
\usepackage{xcolor,colortbl}
\usepackage{url}
\usepackage[round]{natbib}
\usepackage{booktabs}
\newcommand{\eng}[1]{\textit{#1}}
\newcommand{\cmn}[1]{\textit{#1}}
\newcommand{\fre}[1]{\textit{#1}}
\newcommand{\jpn}[1]{\textit{#1}}
\newcommand{\msa}[1]{\textit{#1}}
\newcommand{\txx}[1]{\textbf{#1}}
\newcommand{\lex}[1]{\textbf{\textit{#1}}}
\newcommand{\iz}[1]{\textbf{\texttt{#1}}}
\usepackage{url}
\newcommand{\fd}[1]{\url{#1}}  %% SQL field
\newcommand{\tg}[1]{\url{#1}}  %% POS tag
\newcommand{\wn}[3]{\path{#1}$_{#2_#3}$}
\newcommand{\of}[1]{\textbf{\texttt{#1}}}
\newcommand{\mc}{\multicolumn}
\newcommand{\furl}[1]{\footnote{\url{#1}}}
\usepackage[normalem]{ulem}
\newcommand{\ul}{\uline}
\newcommand{\ull}{\uuline}
\newcommand{\uwl}{\uwave}
%\newcommand{\wn}[3]{\path{#1}\ensuremath{^{#2}_{#3}}}

\usepackage{CJKutf8}
\newcommand{\ja}[1]{\CJKfamily{min}#1\CJKfamily{gbsn}}
\newcommand{\zh}[1]{\CJKfamily{gbsn}#1\CJKfamily{min}}
\usepackage{mygb4e}
\usepackage{pst-node}
\usepackage{graphicx}
\DeclareUnicodeCharacter{2A00}{\ensuremath{\odot}} %⨀
\DeclareUnicodeCharacter{2A01}{\ensuremath{\oplus}} % ⨁



\title{NTU Multi-Lingual Corpus Manual
\\ \huge ongoing notes --- do not cite}
\author{Francis \textbf{Bond}
  \\ Liling \textbf{Tan} and \textbf{Wang} Shan}
% \date{\today}

% % 10 pages max
% % single line spacing
\begin{document}
\begin{CJK}{UTF8}{gbsn}
% 勉強
% \ja{勉強}
\maketitle
\tableofcontents{}
\chapter{NTU Multi-lingal Corpus Motivation}

% \begin{abstract} % no more than 300 words
% % In no more than 300 words, please provide a succinct and accurate
% % description of the proposal to include the specific aims, hypotheses,
% % methodology and approach of the research proposal, including its
% % academic significance \textbf{and relevance to Singapore}.
%   We propose to study how meaning is expressed differently in different
%   languages.
%   This problem has been extensively discussed in translation studies,
%   where it is a major source of translation shifts:
%  ``departures from formal correspondence in the process of going
%   from [the source language] to [the target language]''.
%   However, large empirical studies of meaning differences are rare
%   because meaning is not marked explicitly in text, and there are few
%   corpora with semantic annotation, especially for Asian languages.
%   We will conduct a qualitative and quantitative study of translation
%   shifts on parallel texts in four typologically diverse languages: English, Mandarin
%   Chinese, Indonesian and Japanese (produced by previous research at NTU).
%   The key of our proposal is to make full use of a rich multi-lingual
%   semantic lexicon (Wordnet, which we have extended to cover Japanese
%   and Malay/Indonesian) and use it to link the same meanings across the different
%   languages.
%   In parallel sentences, translation equivalents which do not share
%   the same semantic class are then evidence of translation shifts.
%   Using the multilingual network of
%   semantic classes, we can measure explicitly differences in meaning,
%   even across languages.
%   This will thus be the first large scale quantitative study of
%   translation shifts for Asian languages.
%   Our preliminary study linking semantic classes in Chinese and
%   English found many translation shifts to be caused by predictable differences in
%   syntactic categories and language-specific multi-word expressions.
%   To make the results most widely useful, we will release all of the
%   cross-lingually annotated data along with the extended multi-lingual lexicon.
%   The findings of this study are of broad interest to Singapore,
%   as English, Chinese and Malay are the most commonly spoken official 
%   languages of Singapore, and Japanese is an important regional
%   language.   Further, there are immediate applications to natural
%   language processing tasks such as machine translation and sentiment
%   analysis.
% \end{abstract}

% TODO:
% \begin{itemize}
% \item FIXME: replace the definition of translation shifts by another
%   definition, possibly informal?
% \item FIXME: most of the abstract is very dense
% \item emphasize the significance of the project socially
% \item emphasize the significant contribution of the project to
%   cognitive linguistics:
%   \begin{itemize}
%   \item we are comparing 4 (5?) typologically different languages,
%   \item the potential to capture the conceptual network that underpins
%     these languages will be extremely valuable to the study of human
%     cognition.
%   \end{itemize}
% \end{itemize}

% \newpage
% Original contrib by FB:
%
% It is well known that different languages present the same information
% in different ways.
% These differences can be clearly seen in translation, where they are
% studied in translation theory and corpus linguistics as Translation
% Shifts \citep{Bakker:Koster:Leuven-Zwart:2009,pado2010translation} and
% in machine translation as Translation Divergences \citep{Dorr:1993}.
% However, there are almost no quantitative studies of translation
% divergences due to the expense of annotation.
% We will use advances in WSD to automatically mark divergences in
% several languages, on a scale as not yet seen.
%

The question of how to represent meaning is perhaps the most important
subject in current linguistic research.  It has profound implications
for studies of human cognition, as well as studies of syntactic
structure.  It is also important for natural language processing,
knowledge representation being the core technology underlying the
semantic web.  This proposal directly addresses how to represent
meaning, using differences in language to shed light on exactly what
information is stored and how.
% Translation shifts
It is well known that different languages present the same information
in different ways.  These differences can be clearly seen in
translation, where they are studied in translation theory and corpus
linguistics as Translation Shifts
\citep{Bakker:Koster:Leuven-Zwart:2009,Pado:Erk:2010} and in
machine translation as Translation Divergences \citep{Dorr:1993}.
\citep[p.~73]{catford1965linguistic} originally defined translation
shifts as ``departures from formal correspondence in the process of
going from [source language] to [target language]''.  Such differences
are of great scientific interest, as they allow us to explore how
human languages vary.  Mapping the different ways in which different
language groups represent meaning has the potential to outline the
conceptual network that underpins human cognition.  Classification of
translation shifts is also of practical interest, providing insights
for language learners and translators (both human and machine).

% description
Translation shifts can be broadly divided into two categories
\citep{Pado:Erk:2010}: \emph{grammatical shifts} and
\emph{semantic shifts}.  A grammatical shift consists in the
expression of the same concept, in different languages, by words with
different morphosyntactic properties, such as syntactic category or
number.  For example, there is a grammatical shift in the translation
of \eng{in silence} when it is translated
as 一声\,不响 \cmn{y\={\i}sh\={e}ng\,b\`{u}xi\`{a}ng} ``keeping
silent'': the English prepositional phrase is translated into Chinese
as a verb plus adjective.  A semantic shift is characterized by a
variation, across languages, of the explicit semantic information that
is used to convey the same general meaning.  Across languages, the
same state of events can be described at different levels of precision
(hypernyms), from different points of view (\eng{raise}
vs. \fre{atteindre} ``rise''), or with different numbers of participants
overtly expressed.  An example of this would be \eng{lad} in English
translated as 孩子 \cmn{h\'{a}izi} ``child'': the meaning is similar,
but not identical. In the most extreme cases, equivalent sentences do
not describe the same event but distinct, pragmatically related,
events.  %%[EXAMPLE]

% origin
Some examples of translation shifts are:
\begin{itemize}
\item lexical gaps (some languages make finer meaning distinctions than others for some concepts)
e.g.,  \eng{river} (Eng) vs. \fre{rivi\`{e}re} ``river which flows to the sea'' or \fre{fleuve} ``river which flows into another river'' (Fre).  These are often culturally based: most Asian languages (including Chinese, Malay and Japanese) have different terms for rice as uncooked grains and cooked rice. (sem)
\item part-of-speech gaps (different languages use different syntactic categories to describe the same phenomenon)
ex: adjectives such as \eng{wet, dry, \ldots} are only expressible by stative
verbs in Japanese: \ja{濡れて\,いる} \jpn{nurete\,iru} ``being wetted'' (gram)
\item different argument structures ex: \ja{歩く} \jpn{aruku} ``walk'' has
no goal in Japanese, you cannot ``walk to somewhere'', instead you
must ``go somewhere walking '' (gram)
\item adaption (sometimes a whole concept is replaced by a cultural
  equivalent) \eng{all ones Christmases come at once} (Eng) translated
  as \ja{盆\,と\,正月\,が\,一緒に\,来た\,よう} \jpn{shougatu to bon ga isshoni
    kita you} ``as though summer holiday and New Year come together''
  (sem)
% \item and other more or less cultural language-specific preferences,
% e.g. prefer the use of a more common lexical item or 
% make explicit some information.  (sem)

\end{itemize}



% Empirical studies
Empirical studies of translation shifts are mainly conducted on
parallel corpora, which are collections of texts accompanied with
their translation in one or more languages.  
% (Lack of) Quantitative studies
There are almost no quantitative studies of translation divergences,
mostly due to the expense of annotation.  Another reason for the lack
of quantitative studies is that, historically, much of the study on
translation divergences in machine translation was done when making
rules for the translation lexicons \citep{Dorr:1993}.  Although
motivated by examples in texts, the rules were rarely linked to
examples, and typically do not have any associated information on
their frequency distribution (i.e. how often each rule is applicable
in a given corpus).

% Use of word sense annotation
We are able to study these phenomena quantitatively now due to recent
advances in the availability of lexical resources.  The PI has helped
create and extend semantic lexicons for Chinese, Japanese and
Malay/Indonesian, all linked to the English Wordnet, a well known
semantic network of English
\citep{Xu:Gao:Qu:Huang:2008,Isahara:Bond:Uchimoto:Utiyama:Kanzaki:2008,_Fellbaum:1998}
and Nurril Hirfana et al. (2011).\nocite{Noor:Sapuan:Bond:2011} This
gives us a single combined network, where a concept such as
\textsc{hand$_1$} ``the (prehensile) extremity of the superior limb''
is linked to representations in multiple languages: \eng{hand} (Eng),
手 \cmn{sh\v{o}u} (Cmn), 手 \jpn{te} (Jpn), \msa{tangan} (Mly, Ind),
\fre{main} (Fre), \ldots \citep{Bond:Paik:2012}.  A different concept,
like the \textsc{hand$_2$} in a game of cards, is a separate node, and
will have a different set of representations.  The nodes are linked to
each other with various semantic relations, such as hypernymy (a
\textsc{hand$_1$} is a \textsc{body part$_1$}), meronomy (a
\textsc{finger$_1$} is part of a \textsc{hand$_1$}), derivation (the
verb \textsc{hand$_a$} ``to hand someone something'' is derived from
the noun \textsc{hand$_1$}) and many more.

Further research by the PI has led to several corpora, all with
multiple translations, being annotated with these concepts in the NTU
Multi-lingual Corpus \citep{Tan:Bond:2011}.  Currently there are around
6,000 sentences annotated for sense in Chinese, English and Japanese,
with 2,000 of them also annotated in Indonesian.  They come from four
genres: essays, short stories, advertising copy and news text.  These
detailed annotations of meanings allow us to measure semantic
differences accurately.  Preliminary studies on the sense annotation
of parallel texts in English and Chinese indicate a
certain systematicity in the translation shifts between the concrete
form that these languages typically use to express the same concept.
Consequently, we hypothesize that sense annotation provides accurate
evidence of translation shifts and provides insight into their
characterization.

% Making it feasible: Automation
% We will use advances in WSD to automatically mark divergences in
% several languages, on a scale as not yet seen. 

% % Interpretative hypothesis
% We additionally expect part-of-speech gaps to delineate or follow the
% boundaries of semantic classes, as in the Japanese adjectives example.

This research builds on earlier work done as part of the JSPS-NUS/NTU
Joint Research Project \textit{Revealing Meaning Using Multiple
  Languages} (2010-2012), in which we did the monolingual corpus
annotation, built the wordnet Bahasa and enhanced the Chinese and
Japanese wordnets.  We have currently also applied for a 36,650 Euro
STIC-Asie grant in France on \textit{Deep Integration of Multi-lingual
  Wordnets}, where we would work with partners in Europe and Asia to
improve the quality of the combined wordnets.\footnote{This is with
  partners in INRIA Alpage (PI: France), Linton (Malaysia), BPPT
  (Indonesia), Ljubljana and JSI (Slovenia).}  We further intend
to continue our multilingual analysis by extending it more to syntax
and structural semantics.  We are developing this theme in a Tier Two
grant application on an integrated multi-lingual lexical and
structural representation of meaning.\footnote{This is with partners
  at computer science at NTU; CSLI, Stanford (USA), Oslo (Norway) and
  Kyung-Hee University (Korea).}


\section{Specific aims}
The general aim of this project is to conduct a corpus-based study of
differences in the expression of meaning across languages, taking
as evidence the translation shifts found in the NTU Multi-lingual
Corpus (NTU-MC) and characterizing these differences as patterns of
structural mismatches in the multi-lingual semantic lexicon Wordnet.
This involves the following concrete aims:
\begin{enumerate}
\item Produce and release a new version of the NTU-MC that contains
  links between the word sense annotations across
  languages;
\item Conduct a quantitative and qualitative analysis of the word
  sense mismatches of this corpus, where the qualitative part is
  grounded in the structure of the multi-lingual semantic lexicon
  Wordnet: in this research the links found above will be classified
  into different types;
\item Improve the precision, coverage and density of the multi-lingual
  Wordnet.
\end{enumerate}


\section{Significance}
% in general
How and why the same meaning is conveyed differently across languages
is the foundational question of both contrastive linguistics and
translation studies.
Although their theoretical perspectives and concrete aims sensibly
differ, corpus-based studies conducted in these two disciplines rely
on similar methodologies, tools and resources, among which parallel
corpora occupy a central place~\citep{granger2003corpus}.
If quantitative studies on lexical occurrences and collocations are
routinely performed on parallel corpora nowadays, they still prove
hard to conduct on syntactic patterns and even more so on semantic
categories~\citep{ebeling1998contrastive, granger2003corpus}.

The same foundational question of how and why different languages use
different means to convey the same meaning poses an equally difficult
challenge to concrete applications of natural language processing such
as machine translation and sentiment analysis.
% [FIXME] \cite{Dorr:1993} on issues with translation shifts for machine
% translation
An analysis of the range and distribution of divergences can guide
machine translation researchers as to what kind of transfer rules are
necessary.  In addition, the sense annotated parallel texts and the
enriched semantic lexicon allow (machine) translation rules to be
inferred directly, as well as paraphrase and entailment rules
(which are useful for measuring sentence similarity).

% Related work
Over the past few years, several research projects have been
annotating parallel corpora in order to enable deep corpus-based
studies of translation shifts and cross-lingual differences.

% FuSe
Among the first major efforts in this direction was
FuSe~\citep{cyrus2006building}, an English-German parallel corpus
extracted from the EUROPARL corpus of parliamentary proceedings from
the European Parliament~\citep{Koehn:2005}.
The parallel sentences are first annotated monolingually.
Each predicate is tagged with its lemma and part of speech;
related predicates, e.g. a verb and its nominalization, are marked as
belonging to the same group.
Arguments are tagged with names of pseudo-roles derived from the verb
e.g. \textit{buyer} for \textit{buy}.
Translation shifts are then annotated explicitly to distinguish
between 12 different types, 6 for grammatical shifts and 6 for
semantic shifts.
This project presents two major drawbacks.
First, its annotation scheme is idiosyncratic and is not backed up by
any semantic resource, be it monolingual or cross-lingual, like
Wordnet or FrameNet~\citep{baker1998berkeley}.
Second, the resource has not been released and is not developed
anymore.

% LinES
LinES~\citep{ahrenberg2007lines} is a parallel treebank of $2,400$
sentences from a user manual and a novel, in English and Swedish, with
word alignments and syntactic dependency structures.
Some dependency relations contain a supplementary semantic
information, e.g. for adverbials and prepositional objects, but the
corpus is clearly oriented towards a syntactic study of translation
shifts between English and Swedish.

% SMULTRON
SMULTRON~\citep{Smultron2010} is a parallel treebank of $2,500$
sentences from different genres: a novel, economy texts from several
sources, a user manual and mountaineering reports.
Most of the corpus is German-English-Swedish parallel text, with
additional texts in French and Spanish.
The treebank contains syntactic phrase trees in each language and word
and phrase alignments across languages.
In addition, a preliminary study has been conducted on $50$ sentences
of the English-Swedish parallel subcorpus where each sentence has been
annotated with semantic frames.
The aim was to measure the feasibility of frame projection across
languages, given a syntactic alignment~\citep{Volk:Samuelsson:2007}.
Most of the resource development remains however focused on the
alignment of syntactic structures.

% CroCo
CroCo~\citep{vculo2008empirical} is a German-English parallel and
comparable corpus of a dozen texts from eight genres, totalizing
approximately $1,000,000$ words.
Each sentence is annotated with phrase structures and grammatical
functions, and words, chunks and phrases are aligned across parallel
sentences.
This resource is limited to two languages, English and German, and is
not systematically linked to any semantic resource.
% Moreover, neither the corpus nor its documentation are readily
% available.

% Pado Erk
Last, \cite{Pado:Erk:2010} have conducted a study of translation
shifts on a German-English parallel corpus of $1,000$ sentences from
EUROPARL annotated with semantic frames from FrameNet and word
alignments.
Their aim was to measure the feasibility of frame annotation
projection across languages.
A major difference between their study and our proposal is that frames
have a coarser granularity than Wordnet senses.
For example, verbs that describe different points of view on the same
state of events evoke the same semantic frame.

% Wordnet for contrastive studies
\citet{peters2002cross} showed that multi-lingual Wordnets used in
conjunction with parallel corpora provide a useful basis to conduct
corpus-based contrastive studies.  Such studies are however extremely
rare, \citet{bentivogli2005exploiting} being the main exception,
partly because most Wordnets lack sufficient cover and density.
Density is central to the study of grammatical shifts, as they often
involve replacing the literal translation of a word with another word
that is morphologically related to the former, such as deverbal
adjectives and nouns.  Derivational relations are relatively new in
Wordnet~\citep{fellbaum2003morphosemantic}, therefore their
availability is still scarce~\citep{pala2007derivational,
  fellbaum2009putting, piasecki2010toward, gabor2012boosting}.  As we
investigate and annotate the corpora, we will also increase the
precision, coverage and density of the English, Mandarin Chinese,
Indonesian and Japanese components of the multi-lingual Wordnet.

% Asian languages
All the aforementioned corpus-based studies on translation shifts have
been performed on European languages.  Our project will be the first
corpus-based study of translation shifts involving Asian languages.
We hypothesize that the lexical information and specific preferences
of languages that we will characterize in our study can be applied
onto similar languages of the same family, to test hypotheses or to
increase the coverage of their Wordnets.

% Conclusion of "direct" significance

In this project, we propose a novel approach to quantitative and
qualitative studies of translation shifts on parallel corpora by
making decisive use of a multi-lingual semantic lexicon.  We will
produce and release a parallel corpus with cross-lingual
sense-annotation, which outnumbers the current standards for
comparable experiments~\citep{Pado:Erk:2010} both in scale
(approximately $6,000$ sentences), and diversity, with four languages,
English, Mandarin Chinese, Indonesian and Japanese, and four genres
(essays, short stories, advertising copy and news text).  The
availability of this resource will open the way to further contrastive
and translation studies.  The wordnets are already extremely widely
used, with over a thousand downloads a month, and almost a thousand
on-line queries a day.  The corpora will be directly linked to the
wordnets, and we expect the same level of use.  The findings of our
pilot study will also be useful to concrete applications in natural
language processing.

% broad scientific significance
More broadly, the findings of this study will make a significant
contribution to cognitive linguistics.
We will compare 4 typologically different languages through the prism
of a common conceptual network, which will be extremely valuable to
the study of human cognition.

% social significance
By pointing and explaining meaning differences across languages,
the findings of this study can contribute to the improvement of
teaching methods for second language acquisition and for translation
and interpretation training, as well as mutual understanding and
dialogue across cultures.
These issues are crucial to Singapore as a multilingual and
multicultural society and as a global cultural and business hub.


\section{Approach}

Our approach will proceed in six phases: (0) automatic cross-lingual tagging
(1) manual cross-lingual tagging, (2) monolingual fix-up, (3) final check of
cross-lingual tags, (4) analysis of translation shifts, (5) write-up of
results and release of corpus.

In the zeroth phase, we will create candidate links between content
words in the two languages using the most current version of the
multi-lingual wordnet and statistical models of word alignment.
% In previous research by the investigator at NTU, the wordnets for
% Chinese \citep{Xu:Gao:Qu:Huang:2008}, English \citep{_Fellbaum:1998},
% Japanese \citep{Isahara:Bond:Uchimoto:Utiyama:Kanzaki:2008} and
% Indonesian (Nurril Hirfana et al., 2011)\nocite{Noor:Sapuan:Bond:2011}
% have been combined into a single multi-lingual lexicon linked through
% English as a common pivot \citep{Bond:Paik:2012}.  
%%% FIXME picture here
In a preliminary investigation based on aligning concepts from the
short story \textit{The Adventure of the Dancing Men} (Conan Doyle)
and its Chinese translation \citep{Mok:2012}, 54.6\% of the final
links were between words from exactly the same semantic class, such as
\cmn{sh\v{o}u} and \eng{hand}, which are both part of the semantic
class 05564590-n ``the (prehensile) extremity of the superior limb''.
A further 7.9\% can be linked by allowing a word to link to a word
that is a member of its immediate subclass, such as \eng{forgive}
linking to the more specific: \cmn{r\'{a}osh\`{u}} ``pardon''.  The
remaining 37.5\% of links are less direct (such as non-immediate
hypernyms, siblings, derivationally related words and so forth).  We
will attempt to link more automatically by allowing linking of these
less direct relationships using not just wordnet, but also word-level
alignments using shallow statistical models.  This automatic alignment
will be done by the research fellow, based on the data from the
preliminary studies.

In the first phase, human annotators will check the automatic
alignments and examine any non-aligned words.  Manual alignment is
necessary for the most extreme translation shifts.  All manual
annotation will be carried out by multi-lingual students from LMS, who
have received some training in semantics, under the supervision of the
research fellow.  All data is stored in SQL databases, and annotation
done using web-based tools developed at LMS (and maintained and
extended as necessary by the research fellow).  We will align Chinese,
Japanese, and Indonesian to English for a total of approximately 12,800 sentence
pairs (4,800 Chinese-English (1,200 have been done in preliminary
work), 6,000 Japanese-English and 2,000 Indonesian-English).  Based on
our preliminary studies, we expect this to take around 1,280
person-hours.  To do this we will assemble a team of 4--8
undergraduate annotators working over the summer.

%%% How much will we align
%% C-E 6000 (-1200) = 4800
%% M-E 2000
%% J-E 6000
%% Total (+ 8000 4800) 12800 (10/hour)
%% (* 8 1280) 10,240 (/ 12800 40) 320

Typically, a close examination of the cross-lingual meaning reveals
some errors and omissions in the original annotation.  These will be
corrected in phase (2).  In the preliminary experiments, new word
meanings needed to be created for around 1.6\% of the English words
and 11.6\% of the Chinese words (the non-English wordnets are less
fully developed than the English wordnet).  To make our subsequent
analysis be about the real differences, not just differences in the
lexical resources for each language, it is necessary to fix these
before the final analysis.  It is important to note that we expect to
find fewer missing words in each corpus we annotate. Because word frequency
is skewed, with a few common words appearing very often, once they
have been added the number of omissions soon drops.  New candidate entries
will be made by the student annotators and then vetted by the PI and each
wordnet's development team before being added.  We expect to add
1--2,000 new entries to the wordnets: each entry needs to have a
definition written, and be linked by one or more semantic relations to
the existing structure.

%% say 5/hour: (* 8 (/ 2000 5) 400) 3200

In phase (3), we will review the cross-lingual annotation, with the
monolingual tags corrected.  This should be relatively fast, and will
require fewer annotators.

%% Assume 5 times faster: 2,000
% (/ 1280 5) 256
%(+ 1280 400  256) 1936

Phase (4) is the most important step: here we will analyze the
cross-lingual tags, using the information in the wordnets to model the
differences systematically.  Many kinds of semantic shift will be
identifiable automatically: especially simple syntactic shifts (as
evidenced by different POS-tags) and semantic translation shifts such
as hyponymy/hypernymy, equivalence and adaptation.  Others will
require further linguistic analysis.  Phase (4) will be the main
responsibility of the research fellow.  

Finally, in the final phase (5), we will disseminate the results in
two ways: the first is as academic papers, both in journals and
specialist conferences (such as the Global Wordnet Conference).  The
second is by the release of the enhanced wordnets (in cooperation with
their developers) and the annotated multi-lingual corpora \citep[as
part of the NTU Multi-lingual Corpus:][]{Tan:Bond:2011}.  The corpora
will also be used to add word sense frequency information to the wordnets.


% Distinguish in the annotation between:
% - syntactic/structural translation shifts, as evidenced by different POS-tags, head switching\ldots
% - semantic translation shifts, distinguished between hyponymy/hypernymy, equivalence, adaptation\ldots

We give a more detailed example below.  Consider the sentence \eng{The
  minister left the party}, with translations as given below:
\begin{exe}
  \ex \glll Jpn: \rnode{AJ}{\ja{大臣}$_1$} \ja{が}  \rnode{BJ}{\ja{離党}} した  \\
  {} daijin ga ritou shita  \\
  {} minister \textsc{sbj} leave-party did \\

  \ex Eng: \eng{The \rnode{AE}{minister$_4$} \rnode{BE}{left$_8$} the \rnode{CE}{party$_1$}}

 \ex \glll Cmn: \rnode{AC}{官员$_1$} \rnode{BC}{离开$_3$} 了 \rnode{CC}{政党$_1$}  \\
  {} guanyuan likai le zhengdang  \\
  {} minister leave already political-party \\

\end{exe}
\nccurve[angleA=-135,angleB=150,ncurv=2]{AJ}{AE}
\nccurve[angleA=-45,angleB=60,ncurv=1.7,linestyle=dashed]{BJ}{BE}
\nccurve[angleA=-90]{AE}{AC}
\nccurve[angleA=-90]{BE}{BC}
\nccurve[angleA=-90,angleB=180]{CE}{CC}

As part of the NTU-MC, all words that are in the wordnets have been
marked with their senses.  English \eng{minister, leave} and
\eng{party} are all multiply ambiguous, but they have been annotated
with their senses meaning ``government minister'', ``step down from''
and ``political party'', respectively.  The Chinese sentence has the
same three concepts and can be automatically linked in phase~0.
Although the concepts are exactly the same, the Chinese is in fact
less ambiguous: \cmn{guanyuan} cannot be a religous minister, and
\cmn{zhengdang} can only mean political party.  Therefore there is a
semantic translation shift involving specialization.  The fact that
these meanings are more specialized in Chinese (and that the English
chose not to use the more specific multi-words \eng{government
  minister} and \eng{political party}) can be automatically discovered
from the structure of wordnet.

Looking at the Japanese, there is one word \jpn{ritou} ``leave the
(political) party'' which does not currently exist in the wordnet.
Manual linking in phase~(1) would link it as a hyponym (sub-class) of
\textsc{leave}$_8$ with \textsc{party}$_1$ linked as a pertainym
(related in some way).  In phase~(2) a proper entry for \jpn{ritou}
would be created.  In phases~(3) and~(4), we anticipate that many
similar examples will appear (it is reasonably common for a verb and
its object in English to map to a single verb in Japanese), and we can
look for patterns of this occurring (e.g., is the verb often a
\textsc{verb of motion}?).

This extremely rich annotation allows us to explicitly measure
translation shifts, and also use computational search to look for
patterns in the shifts.  No one has yet carried out such detailed
cross-lingual annotation on such a large scale: either in terms of the
number of sentences or the number of different language families.  By
creating such data for the languages of Singapore and the surrounding
region, we are laying the foundations for research not only on
translation shifts, but for further research on cognitive semantics,
machine translation, text-mining and many others.

The deliverables will include not only several papers, but also three
software packages: (i) the annotated corpora themselves, (ii) the
enhanced wordnets and (iii) tools for automatically annotating and
then manually correcting more data.  As the most direct extension of
this research, we intend to continue by automatically annotating more
data in more languages, and also to investigate techniques for
learning word meanings in new languages by comparing words with their
translations.

% \section{Investigator}

% Francis Bond (PI) is a computational linguist with extensive
% experience in machine translation and natural language understanding.
% He has published over a hundred refereed papers, and has several
% patents.  In addition to his position as Associate Professor at NTU,
% he is an adviser at the Japanese government's National Institute for
% Information and Communications Technology (NICT).  The PI is the lead
% developer for the Japanese Wordnet, the Wordnet Bahasa and the open
% multilingual Wordnet (with over twenty languages represented).  He
% will be responsible for the overall coordination and design of the
% annotation framework as well as the wordnet development.

% Mathieu Morey (research fellow) is a computational linguist
% specializing in efficient representation of linguistic models as
% graphs, both for syntax and semantics.  The RF would be funded by this
% grant, and his role would be to identify and analyze the translation
% divergences automatically.  He would also supervise the cross-lingual
% annotation.


% \section{Environment}

% Singapore's multi-ethnic, multi-lingual society makes it the perfect
% location for research on language differences. The division of
% Linguistics and Multilingual Studies (LMS) at Nanyang Technological
% University is an ideal environment of this project due to its focus on
% empirical, multi-lingual linguistics research.  The research fits
% clearly into two of the University's Five Peaks of Excellence: New
% Silk Road and Innovation Asia.  The computational group at LMS has
% been very successful in involving undergraduates in its research both
% as annotators and as researchers, with several student papers
% published at refereed international conferences.  We think they are
% well suited to investigating the differences in language's
% representation of meaning.  We now have a small, well-annotated
% collection of Chinese-English-Japanese corpora built in cooperation
% with researchers in Japan, to which we are adding Indonesian text.
% For infrastructure, LMS has its own small cluster for computational
% experiments with local servers for annotation.  Outside LMS, there is
% a very active collaboration with other regional wordnet developers (in
% China, Indonesia, Malaysia and Japan) and the Global Wordnet
% Association.  Our research is interdisciplinary and we cooperate with
% computer science researchers at CSE in NTU and at NUS in setting up a
% general framework for representing semantics computationally.
% Finally,  NTU is a member of the Erasmus Mundus Multi exchange program
% and is benefiting greatly from exchanges with groups in
% Europe. Mathieu Morey is currently funded as a post doctoral research
% fellow under this exchange program.

% Data:
% - NTU-MC
% - multi-lingual WN (+ STIC-Asie)

% Infrastructure:
% - DB / framework FYP NUS (should be well advanced)
% - multi-lingual WN



\section{Preliminary studies and progress reports}

We have already worked extensively on many of the components of this
research: developing the wordnets for Chinese, Japanese and Malay/Indonesian as
well as linking them to English
(\citealp{Isahara:Bond:Uchimoto:Utiyama:Kanzaki:2008,Bond:Isahara:Fujita:Uchimoto:Kuribayashi:Kanzaki:2009,Bond:Isahara:Kanzaki:Uchimoto:2008};  Nurril Hirfana et al., 2011)\nocite{Noor:Sapuan:Bond:2011},
annotating the parallel texts monolingually
\citep{Kuroda:Kuribayashi:Bond:Kanzaki:Uchimoto:2011,Tan:Bond:2011,Bond:Baldwin:Fothergill:Uchimoto:2012},
and linking meanings across languages
\citep{Bond:2010,Frermann:Bond:2012}.


Over the summer, two undergraduate interns (one from St John's
College, Santa Fe visiting NTU in Singapore, one from NTU visiting
NICT in Japan) completely annotated one small corpus for word class
correspondences between Chinese and English, and started with another.
This has led to our initial estimates of the numbers of translation
shifts, as well as a submitted paper on differences in the use of
idioms in translation between English and Chinese
\citep{Mok:2012,Kng:Bond:2012}.

\section{Sense Tagging}
\label{sec:sense-tagging}

One goal is to build systems capable of tagging new text: so we do not
just want to fix errors in the corpus but also consider how the tagger
can be improved so new text is treated correctly.


\subsection{Monolingual Sense Annotation}\footnote{Some of this has
  been discussed in \citet{Bond:Wang:Gao:Mok:Tan:2013}}.



First, the Chinese, Japanese and Indonesian corpora were automatically
tokenized and tagged with parts-of-speech.  Secondly, concepts were
tagged with candidate synsets, with multiword expressions allowing a
skip of up to 3 words.  Any match with a wordnet entry was considered
a potential concept.

These were then shown to annotators to either select the appropriate
synset, or point out a problem. The interface for doing sense
annotation is shown in Figure~\ref{fig:tagging}.

\begin{figure*}[!tpb]
  \centering
\includegraphics[width=0.9\textwidth]{wnxling-img1.eps}  
   \caption{Tagging the sense of \textit{cane}.}
   \label{fig:tagging}
\end{figure*}


\subsection{Error Tags}
\subsubsection{Phase 1}

In Figure~\ref{fig:tagging}, the concepts to be annotated are shown as
red and underlined. When clicking on a concept, its WordNet senses
appear to the right of a screen.  The annotator chooses between these
senses or a number of meta-tags: \textbf{e, s, m, p, u}.  Their meaning
is explained below.

\begin{itemize}
\item [\textbf{e}] error in tokenization
  \\ \ul{今} 日　 should be 今日
  \\ \textit{\ul{three-toed}} should be \textit{three - toed}
\item [\textbf{s}] missing sense (not in wordnet) 
  \\ \textit{I program in \ul{python}} ``the computer language''
  \\ COMMENT: add link to existing synset
  \\ $<$\of{06898352-n} ``programming language''
\item [\textbf{m}] bad multiword \\ (i) if the lemma is a multiword,
  this tag means it is not appropriate \\ (ii) if the lemma is
  single-word, this tag means it should be part of a multiword
\item [\textbf{p}] POS that should not be tagged (article, modal,
  preposition, \ldots{})
\item [\textbf{u}] lemma not in wordnet but POS open class (tagged
  automatically)
  \\ COMMENT: add or link to existing synset
\end{itemize}

\subsubsection{Phase 2}

We found that annotators (especially students) had trouble with
distinguishing the error tags.  %%% FIXME Numbers?

We therefore reduced them to three kinds: 


\begin{itemize}
\item [\textbf{e}] error in tokenization/lemmatization: 
  the corpus should be fixed
  \\ \ul{今} 日　 should be 今日
  \\ \textit{\ul{three-toed}} should be \textit{three - toed}
\item [\textbf{w}] wordnet should be enhanced
  \begin{itemize}
  \item Missing synset
    \\ \textit{I program in \ul{python}} ``the computer language''
    \\ COMMENT: add link to existing synset
    \\ $<$\of{06898352-n} ``programming language''
  \item Missing lemma (realization missing for an existing synset)
    \\ COMMENT: add or link to existing synset
  \end{itemize}
\item [\textbf{x}] this lexical unit should not be in wordnet
  \begin{itemize}
  \item bad multiword \\ (i) if the lemma is a multiword,
    this tag means it is not appropriate \\ (ii) if the lemma is
    single-word, this tag means it should be part of a multiword
  \item  POS that should not be tagged (article, modal,
    preposition, \ldots{})
  \end{itemize}
\end{itemize}


\subsection{Missing Senses}

Missing senses in the wordnets were a major issue when tagging,
especially for Chinese and Japanese.  We allowed the annotators to add
candidate new senses in the comments; but these were not made
immediately available in the tagging interface. As almost a third of
the senses were missing in Chinese and Japanese, this slowed the
annotators down considerably.

Our guidelines for adding new concepts or linking words to existing
cover four cases:
\begin{itemize}
\item [=] When a word is a synonym of an existing word, add
=synset to the comment: e.g.\ for laidback, it is a synonym of
\of{02408011-a} ``laid-back, mellow'', so we
add \texttt{=02408011-a} to the comment for laidback.
\item[\textless]  When a word
is a hyponym/instance of an existing word, mark it with
$<$synset: For example, \textit{python} is a hyponym of
\of{06898352-n} \lex{programming language}, so we
add \url{<06898352-n} to \textit{python}
\item[!] Mark antonyms with !synset.
\item[$\sim$] If you cannot come up with a more specific
relationship, just say the  word is related in some way to an existing
synset with  $\sim$synset; and add more detail in the comment.
\end{itemize}


Finally, we have added more options for the annotators: \textbf{prn}
(pronouns) and seven kinds of named entities: \textbf{org}
(organization); \textbf{loc} (location); \textbf{per} (person);
\textbf{dat} (date/time); \textbf{num} (number); \textbf{oth} (other)
and the super type \textbf{nam} (name).  These basically follow
\citet[p207]{Landes:Leacock:Fellbaum:1998}, with the addition of
number, date/time and name.  Name is used when automatically tagging,
it should be specialized later, but is useful to have when aligning.
Pronouns include both personal and indefinite-pronouns.  Pronouns are
not linked to their monolingual antecedents, just made available for
cross-lingual linking.
\chapter{The Sub Corpora}


\section{Tourism (Your Singapore)}

\begin{itemize}
\item English
  \begin{itemize}
  \item Location: \url{cl:/home/bond/work/ntu-mc}
  \item Format: new
  \item ToDo:  para, new
  \item Tagged: 43,973 (4,220 to go; 10,198 unknown)
  \end{itemize}
\item Chinese
  \begin{itemize}
  \item Location: \url{cl:/home/bond/work/ntu-mc}
  \item Format: new
  \item ToDo:  para, new
  \item Tagged: 38,679 (1801 to go; 952 unknown)
  \end{itemize}
\item Indonesian
  \begin{itemize}
  \item Location: \url{cl:/home/bond/public_html/cgi-bin/tagi}
  \item Format: old, one file
  \item ToDo: update, para, new
  \item Tagged: 27159 / 41774 (untagged from sentence 1300 to 2196 plus scattered)
    % (* 100.0 (/ 27159 41774.0)) 65.01412361756115 done
  \item  Mark morphological thingies with ⨀⨁ (U2A00, U2A01) (ter, di, nya)
 \\ harimaunya -> harimau⨀ ⨁nya 
\\ Note: English superlatives should also be reduced to base form
  \end{itemize}
\item Japanese: nothing done 
\end{itemize}


\section{Story}

We chose \textit{The Adventure of the Dancing Men} (danc) and
\textit{The Adventure of the Speckled Band} (spec)\footnote{We use the
  standard abbreviations for Sherlock Holmes Stories proposed by
  \citet{Christ:1947}.}  as they are both highly rated stories with
multiple translations into Japanese on the Aozora Bunko site.  The
English is modern enough to be understood, yet the stories (and often
translations) were produced long enough ago to be out of copyright.


\begin{itemize}
\item English: (dm new, sb old)
  \begin{itemize}
  \item Location: \url{cl:/home/bond/public_html/cgi-bin/tagx}
  \item Format: new
  \item ToDo:  para, new (sb: update, consolidate)
  \item Tagged: ???
  \item Source: \url{http://tomoki.tea-nifty.com/tomokilog/2005/12/sherlock_holmes.html}
  \item Source: \url{http://www.my285.com/zt/kenan/bskw/index.htm}
  \end{itemize}
\item Chinese
  \begin{itemize}
  \item Location: \url{cl:/home/bond/public_html/cgi-bin/tagx}
  \item Format: new
  \item ToDo:  para, new (sb: update, consolidate)
  \item Tagged: ???
  \end{itemize}
\item Japanese
  \begin{itemize}
  \item Location: \url{cl:/home/bond/public_html/cgi-bin/xling}
  \item Format: old, separate files
  \item ToDo:  para, new (sb: NOT TAGGED)
  \item Tagged: ???
  \end{itemize}
\item xling
  \begin{itemize}
  \item Location: \url{cl:/home/bond/public_html/cgi-bin/tagx}
  \item dm only
  \end{itemize}
\end{itemize}

\section{Essay}
\begin{itemize}
\item English:
  \begin{itemize}
  \item Location \url{cl.local://home/bond/public_html/cgi-bin/tagx/eng-catb.db}
\\ debugging info \url{cl.local://home/bond/bin/eng-catb.log}
  \item Format: new
  \item ToDo:  finish going through
  \item Tagged: all
  \item Paper: \url{ronf:/home/bond/papers/wn/wn-tag-catb.tex}
  \end{itemize}
\item Chinese
  \begin{itemize}
  \item Location: \url{cl:/home/bond/public_html/cgi-bin/tagx}
  \item Format: new
  \item ToDo:  para, new (sb: update, consolidate)
  \item Tagged: all
  \end{itemize}
\item Japanese
  \begin{itemize}
  \item Location: \url{cl:/home/bond/public_html/cgi-bin/tagx}
  \item Format: new
  \item ToDo:  para, new (sb: update, consolidate)
  \item Tagged: all
  \end{itemize}
\item xling
  \begin{itemize}
  \item Location: \url{cl:/home/bond/public_html/cgi-bin/tagx}
  \item eng-cmn only; incomplete
  \end{itemize}
\end{itemize}

\section{News (kc01, kc02)}
\begin{itemize}
\item English
  \begin{itemize}
  \item Location: \url{cl:/home/bond/public_html/cgi-bin/tag1}
  \item Format: old, separate files
  \item ToDo: consolidate, update, para, new
  \item Tagged: ???
  \item Everything has problems if there is a '[' in the sentence
  \end{itemize}
\item Chinese
  \begin{itemize}
  \item Location: \url{cl:/home/bond/public_html/cgi-bin/tag1}
  \item Format: old, separate files
  \item ToDo: consolidate, update, para, new
  \item Tagged: ???
  \end{itemize}
\item Japanese (kc01 only tagged)
  \begin{itemize}
  \item Location: \url{???}
  \item Format: old, separate files
  \item ToDo: consolidate, update, para, new
  \item Tagged: kc01 only
  \end{itemize}
\item xling
  \begin{itemize}
  \item sentence and bunsetsu links from NICT
  \end{itemize}
\end{itemize}



\section{SemCor}

We will also try to format SemCor in the same way.


\chapter{Database Structure}

Monolingual data and links are kept in separate databases.

To link two monolingual databases and a link, use the SQL ATTACH command.
\begin{verbatim}
sqlite3 eng-jpn-links.db
sqlite> ATTACH  eng-corpus.db as S
sqlite> ATTACH  jpn-corpus.db as T
\end{verbatim}




\section{Monolingual Database}


Databases are organized by genre, with subcorpora within them.  Within
each language, sentence numbers are unique.

Currently there are four genres, \txx{tourism}, \txx{story},
\txx{essay} and \txx{news}.  We may add another \txx{semcor} although
it is itself composed of multiple genres.

\begin{table}
  \centering
\begin{tabular}{rlll}
  Numbers & Genre & Corpus & Comment \\
  \hline
  1      & essay & catb &  \\
  10,000 & story & danc &  \\
  11,000 & story & spec &  \\
  60,000 & news & kc & and links to original numbers \\
  100,000 & tourism & yoursing &  \\
  110,000 & tourism & med & not annotated \\
  120,000 & gloss & wngloss & not annotated  \\
  300,000 & semcor & semcor & not annotated  
\end{tabular}

  \caption{Sentence Numbers in the Corpus}
\end{table}

\section{Monolingual Tables}

\subsection{Meta}

\begin{table}[htpb]
  \centering
  \begin{tabular}{llll}
  \textbf{row}   &  \textbf{type} & \textbf{data} & \textbf{comment} \\
  \hline
  corpus & string & corpus name & name without \texttt{.db}\\
  title & string & display name & \\
  language & string & text language & iso 3 letter  \\
\hline
\end{tabular}
\caption{\texttt{meta}: Meta for words}
\end{table}

\subsection{Subcorpus Info}
\begin{table}[htpb]
  \centering
  \begin{tabular}{llll}
  \textbf{row}   &  \textbf{type} & \textbf{data} & \textbf{comment} \\
  \hline
  scid  & integer & subcorpus & starts from 1\\
  subcorpus & TEXT & short name & same across languages\\
  subtitle &TEXT & display name & \\
  url & TEXT & source & 
\\ \hline
\end{tabular}
\caption{\texttt{subcorp}: Subcorpus info}
\end{table}

\begin{table}[htpb]
  \centering
  \begin{tabular}{llll}
  \textbf{scid}   &  \textbf{subcorpus} & \textbf{subtitle} & \textbf{url} \\
  \hline
  1 & spec & The Adventure of the Speckled Band &  www.??? \\
  2 & danc & The Adventure of the Dancing Men &  www.??? 
\end{tabular}
\caption{\texttt{subcorp}: Subcorpus Example}
\end{table}

% e.g:  1,'sb', 'The Adventure of the Speckled Band', www.??? 

FIXME: maybe add license here?


\subsection{Monolingual Sentence Information}

We keep the original sentence, in case we need to recover from bad
tokenization.  \fd{cfrom} and \fd{cto} are calculated relative to this.

\fd{pid} is currently unused.  We may need more META info (header, \ldots).

\begin{table}[htbp]
  \centering
  \begin{tabular}{llll}
  \textbf{row}   &  \textbf{type} & \textbf{data} & \textbf{comment} \\\hline
  sid   & integer & sentence ID & unique over all corpora\\
  pid   & integer & paragraph ID & unique in subcorpus, start from 0\\
  sent  & text    & sentence & pre-tokenization\\
  scid  & integer & subcorpus & starts from 1\\
  comment & string & comment \\
\multicolumn{4}{l}{FOREIGN KEY(scid) REFERENCES subcorp(scid)}
\\ \hline
\end{tabular}
\caption{\texttt{sent}: Sentence information}
\end{table}

\subsection{Monolingual Word Information}

\begin{table}[htpb]
  \centering\caption{\texttt{sent}: Table for original sentence}
\bigskip
  \begin{tabular}{llll}
  \textbf{row}   &  \textbf{type} & \textbf{data} & \textbf{comment} \\\hline
  sid   & integer & sentence ID & unique over all corpora\\
  wid   & integer & word ID & unique in sentence, start from 0\\
  word  & string  & surface form &  \\
  pos   & string  & part of speech & per language  \\
  lemma & string  & lemma \\
  cfrom$^*$ & integer & cfrom & start character offset of word in sentence \\
  cto$^*$ & integer & cto & end character offset of word in sentence \\
  comment & string & comment \\
\hline
\end{tabular}
\\ $^*$ currently unused
\caption{\texttt{word}: Table for word information}
\end{table}

\subsection{Monolingual Concept Information}

Currently cid is unique in a given corpus (genre).  When we link the
corpora together, this causes problems.

ToDo: make 

\begin{table}[htpb]
  \centering  \begin{tabular}{llll}
  \textbf{row}   &  \textbf{type} & \textbf{data} & \textbf{comment} \\\hline
  cid   & integer & concept ID  & unique in corpus, start from 1 \\ 
  sid   & integer & sentence ID & (sid, wid) gives the word \\
  wid   & integer & word ID & \ldots can have more than one\\
  clemma$^{**}$ & string  & concept lemma & may be multi-word \\
  nlemma$^{**}$ & string  & concept lemma & used for error in corpus \\
  tag   & string & wordnet id, tag, $\ldots$ \\
  tag   & string & wordnet id, tag, $\ldots$ & used for tag not in tags\\
  tags$^{***}$  & string & list of tags & for updating \\
  comment & string & comment \\
\hline
\end{tabular}
\\ $^{**}$ was called lemma
\\ $^{***}$ was called ss
\caption{\texttt{concept}: Table for meaning information}
\end{table}

\newpage
\section{Cross-Lingual Linking Database}

For example: \url{eng-cmn-story-links.db}

\subsection{Cross-Lingual Meta}
\begin{table}[htbp]
\centering
  \begin{tabular}{llll}
  \textbf{row}   &  \textbf{type} & \textbf{data} & \textbf{comment} \\
  \hline
  fcorpus & string & source corpus & name without \texttt{.db}\\
  tcorpus & string & target corpus & name without \texttt{.db}\\
  comment & string & comment \\
\hline
\end{tabular}
\caption{\texttt{meta}: Meta for links}
\end{table}

\subsection{Cross-Lingual Sentence Links}
\begin{table}[htbp]
\centering

  \begin{tabular}{llll}
  \textbf{row}   &  \textbf{type} & \textbf{data} & \textbf{comment} \\\hline
  slid & integer & link ID & unique in corpus pair \\
  fsid   & integer & source sentence ID & \\
  tsid   & integer & target sentence ID & \\
  ltype  & string  & link type &  \verb|=, ~, <, >|, Null if unknown\\
  conf   & float & confidence \\
  comment & string & comment \\
\hline
\end{tabular}
\caption{\texttt{slink}: Table for linking sentences}
\end{table}

\subsection{Cross-Lingual Word Links}

Need to add \fd{fsid}, \fd{tsid} as wid numbers are unique only per sentence.

\begin{table}[htbp]
\centering
  \begin{tabular}{llll}
  \textbf{row}   &  \textbf{type} & \textbf{data} & \textbf{comment} \\\hline
  wlid & integer & link ID & unique in corpus pair \\
  fwid   & integer & source word ID & \\
  twid   & integer & target word ID & \\
  ltype  & string  & link type &  \verb|=, ~, <, >|\\
  conf   & float & confidence \\
  comment & string & comment \\
\hline
\end{tabular}
\caption{\texttt{wlink}: Table for linking words}
\end{table}

\subsection{Cross-Lingual Concept Information}

Need to add \fd{fsid}, \fd{tsid} when cid numbers become unique only per sentence.

\begin{table}[htbp]
\centering
  \begin{tabular}{llll}
  \textbf{row}   &  \textbf{type} & \textbf{data} & \textbf{comment} \\\hline
  clid & integer & link ID & unique in corpus pair \\
  fcid   & integer & source concept ID & \\
  tcid   & integer & target concept ID & \\
  ltype  & string  & link type &  \verb|=, ~, <, >|\\
  conf   & float & confidence \\
  comment & string & comment \\
\hline
\end{tabular}
\caption{\texttt{clink}: Table for linking concepts}
\end{table}




\subsection{Example}

From \url{dm} ``The Adventure of the Dancing Men''.

\includegraphics[width=0.9\textwidth]{graph}

This shows a link that worked: \eng{study}, a link that is possible: \eng{nothing} to \zh{东西}, and a link missed due to a lack of concepts in the Chinese wordnet: \eng{move} to \zh{动}.





% \subsection{Tourism2}

% \subsection{Semcor}

% \subsection{Gloss}

\begin{table*}[tb]
  \centering
  \begin{tabular}{lrrrllllllll}
    Genre & \multicolumn{3}{c}{Size (K)}        & \multicolumn{5}{c}{Languages} & Annotation & Comment \\
          & St & Wd & Cpt &  cmn & eng & jpn & ind & other\\
\hline
    Tourism1 & 3.3$^*$ & 52 & & C &  \textbf{C} & P  & .65C & kor, vie, tha 
    &   & \\
    Story    & 1.2 & & & C &  \textbf{C} & C & --- &  
    & xling  & Other corpora \\
    Essay    & 0.8& & & C & \textbf{C} & C & --- & many 
    & HPSG (eng) \\
    News     & 2.0 & 53.0 & 30.0 & C & C & \textbf{.5C}  & --- &
    & Sense, Dep (jpn) & $10^6$ words\\
    \hline
    Tourism2 & 3.8 & 72  & & P & \textbf{P} & P & P & arb, vie\\
    Semcor & 14.2 &  & 151 & --- & \textbf{C} & 0.38C & ---& & HPSG, Penn\\ % (/ 58.0 151)0.3841059602649007
    Gloss  & 112.0 & & &  --- & \textbf{C} & P & ? & alb, spa, \ldots\\
  \end{tabular}
  \caption{Current State of the Tagged Texts from the NTU-MC}
  \label{tab:ntu-mc-now}
  \begin{flushleft} \small
    \textbf{Bold} is source language
    \\ $^*$ less for non-English $\approx 2$k
    \\ P: POS 
    \\ C: Concepts (WN synsets), implies POS 
  \end{flushleft}
\end{table*}




\begin{tabular}{lrrll}
  Task & Time (h) & Cost & Who & Comments \\ \hline
  cmn new words & 40 &  &     & \\
  ind new words & 40 &  &     & 3,379 'u' or 's' (mainly fw)\\
  jpn new words & 40 & EG & \\
  cmn tagging & 43  & Shan? &     & 4,220 concepts: tool feedback\\\\
  ind tagging & 132 &  &     & 13,125 concepts\\\\
  kc-tagging (auto) & Heng & 160 &  4 weeks \\
  eng-cmn link & 340 &  & & 3,400 sentences (not kc)\\
  eng-jpn link & 340& HTY, EG & & 3,400 sentences\\
  eng-ind link & 220 &  & & 2,197 sentences\\
  checking & \\
\hline
Total & 1,315 \\

\end{tabular}

\begin{itemize}
\item Available Money: 15,488 (1,548 hours at 10/hour)
\item WN tagging: 100 open-class words/hour (had estimated 130 in the past)
\item Linking: 10 sent/hour  (up till now 3/hour!)
\item Final Checking: ?
\item People
  \begin{itemize}
  \item HYT (Yue Hui Ting): 3 weeks (120h) can do Japanese
  \item Heng (Heng Hok Hui): CS student,  knows Japanese, can program
  \item JT (Jeanette): 4th year,knows Japanese, can program (built interface)
  \item Malay student (forgot the name)
  \end{itemize}
\end{itemize}

\subsubsection{Ideas}
\begin{itemize}
\item Japanese tag and project first: link and tag together (yoursing)?
\item Indonesian: project
\item tag monosemous as ok
  \begin{itemize}
  \item reduces Chinese by 1259 (40\%)
  \item reduces Indonesian by 3,119 (23\%) %(/ 3119 13135.0) 0.2374571754853445
  \end{itemize}
\end{itemize}

\subsubsection{Tools}
\begin{itemize}
\item Search/display (add to OMW tool)
  \begin{itemize}
  \item synset look up
  \item word look up
  \item sense look up
  \end{itemize}
\item Tagging
  \begin{itemize}
  \item sequential/targeted tagging
    \begin{itemize}
    \item tag all button
    \end{itemize}
  \item pair tagging (this word in eng and that in jpn)
  \item integration:
    \begin{itemize}
    \item retag a word with the latest wn (auto/manual)
    \item then link it
    \end{itemize}
  \item fix morphology
  \end{itemize}
\item Pre-processing
  \begin{itemize}
  \item mono tagger (done: long tail)
  \item xling tagger (extend)
  \item add paragraphs
  \item mono+xling+suggest new words (later)
  \end{itemize}
\end{itemize}





\section{ToDo}

\begin{itemize}
\item Consolidate corpora (Mar, Apr)
  \begin{itemize}
  \item not everything is in the new format
  \item annotation unfinished: Indonesian, Japanese
    \\ timetable these
  \item tokenization errors: 
    \\ English: mainly '-'
    \\ Chinese: proper names, Chengyu
    \\ Bahasa: ter+adj (how to represent decomposed: two concepts with same wid)
  \item NEs not tagged (try to do at least some)
  \item Easy online lookup
    \\ by word, lemma, POS, synset, hypernym, file
  \item New words (from tags 'u', 's', new wordnets, pron, int, dates)
    \begin{itemize}
    \item Pronouns, question words not tagged:   add by hand
    \item Need tool for new entries
    \end{itemize}
  \end{itemize}
\item Crosslingual annotation (Summer)
  \begin{itemize}
  \item Goal (14,000 sentence pairs)
    \begin{itemize}
    \item C-E 6000 (1200 done)
    \item M-E 2000
    \item J-E 6000
    \end{itemize}
  \item Note: News corpus aligned by NICT: hire student to convert? (4,000!)
  \end{itemize}
\item Papers/Research
  \begin{itemize}
  \item  Comparison of WN cover, estimate of missing words/senses
    \\ in the style of \citet{Baldwin:Bender:Flickinger:Kim:Oepen:2004}
    \\ for the four genre: how much missing, how much overlap
    \\ /home/bond/papers/wn/wn-ml-tag
    \\ ACL short, LREC?
  \item xling annotation paper (based on Gao (FYP) and \citet{Mok:2012})
    \\ PACLIC?, ACL short?
\item Translation Shifts (main research: ACL next year, CL)
\item Transfer Rule Acquisition
\item Xling WSD/Annotation
\item FYP: Pronouns
\item Idioms
  \end{itemize}
\item Misc
  \begin{itemize}
  \item Chinese Wordnet --- release?
    \begin{itemize}
    \item Chinese WN
    \item Gloss corpus
    \item Classifiers
    \item[?]
    \end{itemize}
  \item ZhEn 
  \end{itemize}

\end{itemize}


\subsection{Tourism1}

\begin{itemize}
\item English
  \begin{itemize}
  \item Location: \url{cl:/home/bond/work/ntu-mc}
  \item Format: new
  \item ToDo:  para, new
  \item Tagged: 43,973 (4,220 to go; 10,198 unknown)
  \end{itemize}
\item Chinese
  \begin{itemize}
  \item Location: \url{cl:/home/bond/work/ntu-mc}
  \item Format: new
  \item ToDo:  para, new
  \item Tagged: 38,679 (1801 to go; 952 unknown)
  \end{itemize}
\item Indonesian
  \begin{itemize}
  \item Location: \url{cl:/home/bond/public_html/cgi-bin/tagi}
  \item Format: old, one file
  \item ToDo: update, para, new
  \item Tagged: 27159 / 41774 (untagged from sentence 1300 to 2196 plus scattered)
    % (* 100.0 (/ 27159 41774.0)) 65.01412361756115 done
  \item  Mark morphological thingies with ⨀⨁ (U2A00, U2A01) (ter, di, nya)
 \\ harimaunya -> harimau⨀ ⨁nya 
\\ Note: English superlatives should also be reduced to base form
  \end{itemize}
\item Japanese: nothing done 
\end{itemize}


\chapter{The Wordnets}

\section{Princeton Wordnet}

\section{Chinese Wordnet}

\begin{itemize}
\item ToDo:
  \begin{itemize}
  \item convert NTU$_T$ data
  \item discuss simplified/plain
  \item discuss domains
  \item new words from corpus
  \item other?
  \end{itemize}

\end{itemize}

\section{Wordnet Bahasa}

\subsection{Derivational links}

\section{Japanese Wordnet}


\section{Guide to Extending the Wordnet}

Our goal in building the wordnet is to have coverage over all open
class predicates, whether from single words, multi-word expressions or
constructions.  We do not attempt to cover all proper names (entities)
although we have some common ones

\subsection{What goes in}

\subsubsection{Lemmas (Lexical Units)}

\subsubsection{Pronouns}

See Seah FYP


\subsubsection{Interrogatives}

\subsubsection{Named Entities}

We tag proper named entities with the tags shown in Table~\label{sec:names}.
These follow the tags used in SemCor
\citep[p207]{Landes:Leacock:Fellbaum:1998}, with the addition of
Date/Time.

There doesn't seem to be a perfect supertype for \lex{date}.  We want
something referring to a point in time (like I think \wn{today}{n}{1}
does), but this seems to have period of time as its hypernym.

\begin{table*}
  \centering
  \begin{tabular}{llll}
    \textbf{Tag} & \textbf{Meaning} & Instance-of &\textbf{Example} \\ \hline
    org & Organization (Institution) & \wn{group}{n}{1} & Singapore Airlines \\
    loc & Location (Place)    & \wn{location}{n}{1} & Newton Hawker Centre, Sim Lim Square\\
    per & Person       & \wn{person}{n}{1} &Sir Thomas Stamford Raffles\\
    dat & Date/Time    & \wn{time_period}{n}{1} & 2012\\
    oth & Other        & \wn{entity}{n}{1} & Singapore Food Festival\\
  \end{tabular}
  \caption{Tags for Names}
  \label{tab:names}
\end{table*}

In this case the meta-tag is 'x', and the ntag os e.g. 'org'.

We can use the part of speech tags from some systems to boot strap
these.  For example, the IPA-DIC tags \citep{IPAL} distinguish between
person, organization, location and other names.

In general, words tagged with these will not be put into wordnets,
although wordnet-wikipedia hybrids are beginning to go the other way
and include many of these.

For the yoursing corpus, we added many Singaporean locations and
organizations in four languages (Chinese, Malay/Indonesian, English
and Japanese).

Give some numbers for these.


\subsection{What stays out}

\subsubsection{Exclamatives}
\label{sec:exclamatives}
We follow the Princeton Wordnet Guidelines \citep[pp
???]{_Fellbaum:1998} and do not include interjections and exclamatives.

\begin{verbatim}
# 90000482	eng:lemma	wow
# 90000482	jpn:lemma	ワウ
# 90000482	cmn:lemma	哇
# 90000482	eng:def	(Interjection) Used to express amazement
# 90000482	eng:exe	Wow, that is a beautiful painting!
# 90000482	interjection

# 90000483	eng:lemma	hey
# 90000483	cmn:lemma	嘿
# 90000483	jpn:lemma	おおい
# 90000483	jpn:lemma	ちょっと
# 90000483	eng:def	(Interjection) Used to call for someone’s attention
# 90000483	eng:exe	Hey, stop!
# 90000483	interjection
\end{verbatim}
% \subsubsection{Modals}

\subsubsection{Auxiliaries}

We don't tag auxiliaries (although we should do as a student project)



\chapter{Tagging Projects}

HG202 (2011): Yoursing
HG202 (2012): catb
??? Chinese Class
HG202 (2012): SB

JSPS/NICT
Tier One
Tier Two
URECA


2013 catchup
ntu-mc with new wn: COW+SEW+WIKT+ILI
retagged

duplicate sentences deleted





% \newpage
% \appendix
% \section{General information}
% {\bf Keywords:} translation shift, multilingual processing, wordnet, sense annotation
% \ldots{} (6 max.)\newline
% {\bf Funding requested:} 81,000 \$\newline
% {\bf Duration:} 12 months\newline
% {\bf Proposed dates:} 03/2013 -- 02/2014\newline

% \section{Progress}
% \subsection{Milestones}
% \begin{footnotesize}
%   \begin{tabular}{lp{4cm}cccc}
%     \toprule
%     SN & Description & Year 1 Q1 & Q2 & Q3 & Q4\\
%     \midrule
%     & xxx & x & x & x & x\\
%     \bottomrule
%   \end{tabular}
% \end{footnotesize}


% \subsection{Project deliverables}
% \subsubsection{Manpower Training}
% \begin{footnotesize}
%   \begin{tabular}{lc}
%     \toprule
%     No. of M.Sc. students to be trained & 0\\
%     No. of Ph.D. students to be trained & 0\\
%     No. of Undergraduate/Honours students to be trained & 2\\
%     \bottomrule
%   \end{tabular}
% \end{footnotesize}

% \subsubsection{Targets for Research Outcomes}
% \begin{footnotesize}
%   \begin{tabular}{lc}
%     \toprule
%     No. of Publications & 2\\
%     No. of Conferences & 2\\
%     No. of Patents (if applicable) & 0\\
%     No. of Software Tools (if applicable) & 1\\
%     \bottomrule
%   \end{tabular}
% \end{footnotesize}

% \section{Previous Tier 1 awards}
% \begin{footnotesize}
%   \begin{tabular}{lp{2.2cm}cccccc}
%     \toprule
%     SN & Declaration for & Role & Project Title & Year of Award &
%     Project Start Date & Project End Date & Status\\
%     \midrule
%     \bottomrule
%   \end{tabular}
% \end{footnotesize}

% \section{Declaration of other funding support}
% The Principal Investigator and Co-Investigator(s) should detail the
% grants awarded in the last 5 years, including those currently held or
% submitted for consideration by other funding agencies.
% Please note that parallel submissions are not allowed -
% i.e. applicants should not send similar versions or part(s) of the
% current AcRF proposal application to other agencies for funding.
% \textbf{Please include an abstract for all pending and awarded
%   grants.}

% \subsection{Projects pending approval}
% \begin{footnotesize}
%   \begin{tabular}{lp{1.4cm}cp{1.8cm}p{1.6cm}p{1.5cm}p{1.5cm}p{1cm}p{1cm}p{1.4cm}}
%     \toprule
%     SN & Declaration for & Status & Project Title & Funding Agency &
%     Amount Requested AcRF (\$) & Amount Requested Non-AcRF (\$) &
%     Duration From & Duration To & Proposed Duration of Award (No. of
%     Months)\\
%     \midrule
%     & Francis Bond & co-PI & Int\'{e}gration profonde de wordnets
%     multilingues (Deep Integration of Multilingual Wordnets) & French
%     Ministry of Foreign Affairs & 0 ??? & 72,195.41 & 2012 & 2014 & 24\\
%     \midrule
%     & & & & & 0 & 72,195.41 & & & \\
%     \bottomrule
%   \end{tabular}
% \end{footnotesize}
% % 72,195.41 SGD = 46,650 EUR = 36,650 + 10,000 EUR (Foreign Affairs +
% % Inria in the proposal)
% % 24 months from 10/2012 to 09/2014

% % \vspace{3mm}
% % Question:
% % Should part of the local support we listed in the budget of the STIC
% % Asie proposal, e.g. expected PhD funding, be listed as AcRF funding? NO

% \subsection{Projects awarded}
% \begin{footnotesize}
%   \begin{tabular}{lp{1.4cm}cp{1.8cm}p{1.6cm}p{1.5cm}p{1.5cm}p{1cm}p{1cm}p{1.4cm}}
%     \toprule
%     SN & Declaration for & Status & Project Title & Funding Agency &
%     Amount Requested AcRF (\$) & Amount Requested Non-AcRF (\$) &
%     Duration From & Duration To & Proposed Duration of Award (No. of
%     Months)\\
%     \midrule
%     \bottomrule
%   \end{tabular}
% \end{footnotesize}

% KHU

% \subsection{Projects completed}
% \begin{footnotesize}
%   \begin{tabular}{lp{1.4cm}cp{1.8cm}p{1.6cm}p{1.5cm}p{1.5cm}p{1cm}p{1cm}p{1.4cm}}
%     \toprule
%     SN & Declaration for & Status & Project Title & Funding Agency &
%     Amount Requested AcRF (\$) & Amount Requested Non-AcRF (\$) &
%     Duration From & Duration To & Proposed Duration of Award (No. of
%     Months)\\
%     \midrule
%     \bottomrule
%   \end{tabular}
% \end{footnotesize}

% SUG, JSPS, CC


% \subsection{Other resources available to the team}

% \begin{footnotesize}
%   \begin{tabular}{llll}
%     \toprule
%     Types Of Resources & Funding Organisations & Duration of Support &
%     Expiry Date\\
%     \midrule
%     \bottomrule
%   \end{tabular}
% \end{footnotesize}

% Support from NICT, joint research with NTT

% \section{Details and justification of proposed budget}

% Overall proposed budget:

% \begin{footnotesize}
%   \begin{tabular}{lcccc}
%     \toprule
%     Category & Year 1 (S\$) & Year 2 (S\$) & Year 3 (S\$) & Total (S\$)\\
%     \midrule
%     EOM & 81,000 & & & 81,000 \\
%     Equipment & & & & \\
%     OOE - Materials \& Consumables & & & & \\
%     OOE - Training/Other Miscellaneous & & & & \\
%     OOE - Overseas Travel & & & & \\
%     \midrule
%     Sub-Total (S\$) & 81,000 & & & 81,000 \\
%     \bottomrule
%   \end{tabular}
% \end{footnotesize}

% \subsection{EOM}

% \begin{footnotesize}
%   \begin{tabular}{p{3cm}p{0.5cm}p{0.8cm}p{1.2cm}p{1.2cm}p{1.2cm}p{1cm}p{1cm}p{1cm}p{1cm}p{1cm}}
%     \toprule
%     Staff Category & No. of Staff & Cost per Head (\$) & Cost Type & Year
%     1 Man-month/ Total Hours & Year 2 Man-month/ Total Hours & Year 3
%     Man-month/ Total Hours & Year 1 (\$) & Year 2 (\$) & Year 3 (\$) &
%     Total (\$)\\
%     \midrule
%     Research Fellow & 1 & 66,000 & Salary & 12 M-m & & & 66,000 & & & 66,000\\
%     Student Assistant & 1 & 9,000 & ??? & 1,000 Hrs & & & 9,000 & & & 9,000\\
%     Conference Support & ? & 6,000 & Travel & ? & & & 6,000 & & & 6,000\\
%     \midrule
%     & & & & & & & 81,000 & & & 81,000 \\
%     \bottomrule
%   \end{tabular}
% \end{footnotesize}
% % research fellow: 62,000 to 120,000 p.a.
% % student assistant: 7 to 9/hour for undergrads, 10 to 15 for grads
% % conference support: 30,000 for 3 years (Tier 2), 6,000 for 1 year
% % (Tier 1)

% Should travel money be listed here?
% The format of OOE - Travel makes me hesitate (specify the country).

% \subsection{Equipment}

% \begin{footnotesize}
%   \begin{tabular}{p{1.4cm}p{1.4cm}p{0.8cm}p{0.8cm}p{1cm}p{1cm}p{1cm}p{1cm}p{1cm}p{1cm}p{1cm}p{1.2cm}}
%     \toprule
%     Item & Type & Quantity & Cost per Unit & Year 1 Units & Year 2 Units &
%     Year 3 Units & Year 1 (\$) & Year 2 (\$) & Year 3 (\$) & Total (\$) &
%     Quotation\\
%     \midrule
%     \bottomrule
%   \end{tabular}
% \end{footnotesize}

% \subsection{OOE - Materials \& Consumables}
% \label{sec:ooe-mc}

% \begin{footnotesize}
%   \begin{tabular}{p{1.4cm}cccc}
%     \toprule
%     Item & Year 1 (\$) & Year 2 (\$) & Year 3 (\$) & Total (\$)\\
%     \midrule
%     \bottomrule
%   \end{tabular}
% \end{footnotesize}

% \subsection{OOE - Training/Other Miscellaneous}
% \label{sec:ooe-misc}

% \begin{footnotesize}
%   \begin{tabular}{p{1.4cm}cccc}
%     \toprule
%     Item & Year 1 (\$) & Year 2 (\$) & Year 3 (\$) & Total (\$)\\
%     \midrule
%     \bottomrule
%   \end{tabular}
% \end{footnotesize}

% \subsection{OOE - Overseas Travel}
% \label{sec:ooe-travel}

% \begin{footnotesize}
%   \begin{tabular}{p{1.4cm}ccccc}
%     \toprule
%     Country & No of Days & Year 1 (\$) & Year 2 (\$) & Year 3 (\$) & Total
%     (\$)\\
%     \midrule
%     \bottomrule
%   \end{tabular}
% \end{footnotesize}

% \section{Suggested names of referees}
% \label{sec:referees}

% \begin{footnotesize}
%   \begin{tabular}{cp{2cm}p{2cm}p{2.6cm}p{2cm}p{2cm}p{2.6cm}}
%     \toprule
%     SN & Name & Designation & Institution & Email & Contact Number & Remarks\\
%     \midrule
%     \bottomrule
%   \end{tabular}
% \end{footnotesize}
\section{Tagging}

\subsection{Japanese}
\CJKfamily{min}
We introduce two new parts of speech, to deal with verbal-noun
light-verb combinations like 勉強\ する \jpn[study (lit: study
do)]{benky\=o suru}.  In this case, we want to tag
the 勉強 \jpn{benkyou} as a verb, not a noun, and not
tag する \jpn{suru}, instead treating it as an auxiliary.

\begin{table}
\begin{tabular}{lllll}
  POS1 & POS2 & gloss &Example & Comment\\
  \hline
  名詞 & サ変接続 & n-sahen & \ul{勉強} が　好き \\
  動詞 & サ変接続 & n-sahen & 日本語 を \ul{勉強} する & new\\
  動詞 & 接尾  & v-suffix & 日本語 を 勉強  \ul{する/なさる/致す/いたす} & \\
  　　 &  & & 日本語 が 勉強  \ul{できる} \\
\end{tabular}
\caption{New Parts of Speech for Japanese}
\end{table}
Note that the POS code \jpn[v-suffix]{動詞 接尾} exists in the IPA
tagging guidelines, but is not normally used to mark these light verbs.

We identify them as follows:
\begin{verbatim}
select  a.sid, a.wid, b.wid, a.lemma, b.lemma 
from word as a left join word as b   
where a.sid=b.sid and b.wid = a.wid +1
and b.lemma in ('できる', 'する', 'なさる', '致す', 'いたす') 
and a.pos = '名詞-サ変接続';
\end{verbatim}

FIXME: need to deal with cases like these: currently fix する,  but should also make an entry for ``尖鋭 化''.

\begin{verbatim}
    <wf wid='w3573:13' sent='3573' para='0'>尖鋭</wf> <!-- 名詞-一般 -->
    <wf wid='w3573:14' sent='3573' para='0'>化</wf> <!-- 名詞-接尾-サ変接続 -->
    <wf wid='w3573:15' sent='3573' para='0'>する</wf> <!-- 動詞-接尾 -->
\end{verbatim}

Code in \url{/home/bond/ntu-mc/2013-12-14/fix-jpn.py}

\CJKfamily{gbsn}
\subsection{Indonesian}

We did four kinds of morphological analysis, on top of a tagger trained on FIXME.

\begin{itemize}
\item Splitting \lex{di-} ``passive'' from verbs.  Look for
  \textit{\^di+REST} where \textit{REST} is in wordnet as a verb.
\item Splitting \lex{ter-} ``superlative'' from adjectives.  Look for
  \textit{\^ter+REST} where \textit{REST} is in wordnet as an adjective.
\item Splitting \lex{-nya} ``pronoun clitic'' from nouns.  Look for
  \textit{REST-nya\$} where \textit{REST} is in wordnet as a noun, and the whole thing is not in a list of exceptions: FIXME
\item Splitting \lex{-lah} ``please'' from adjectives, adverbs,
  determinatives and verbs.  Look for \textit{REST-lah\$} where
  \textit{REST} is in wordnet as an adjective, adverb, \lex{ini} or
  \lex{itu} and verbs.
\end{itemize}

\subsection{English}

\paragraph{Operating System} lemmatized as \lex{operate system} so only
matches if we check surface form for matching MWEs.


\lex{DJ} is a verb, \lex{dj} a verb.  Both should be upper case?

\subsection{Chinese}

The lemma is different from the surface form for three
cases:

\subsubsection{Reduplication}
大大 AD2  大

We do this by checking for a repeated character, and if not on a stop
list, then add '2' to the POS and delete one character from the lemma.

FIXME: add stop list
FIXME: table showing frequency

AD2
JJ2        
M2        
NN2        
PN2        
VA2        
VV2       
\subsubsection{Plural}

 We also lemmatized plural-marked nouns, such
as 学生+们 \cmn[student+s]{xuéshēng+men}
to 学生 \cmn[student]{xuéshēng}.  This only occurs for 18
words.\footnote{In some segmentation standards these would be two
  tokens, the Chinese Penn Treebank consistently treats these as one
  token \citep{Xia:2000}.} 

\subsection{Errors}

Old characters.
FIXME: expand, list

\subsection{Inconsistent Morphological analysis:}
\begin{itemize}
\item 1990|年 vs １９９４年
\end{itemize}


\section{ToDo}

\begin{itemize}
\item Consolidate corpora (Mar, Apr)
  \begin{itemize}
  \item not everything is in the new format
  \item annotation unfinished: Indonesian, Japanese
    \\ timetable these
  \item tokenization errors: 
    \\ English: mainly '-'
    \\ Chinese: proper names, Chengyu
    \\ Bahasa: ter+adj (how to represent decomposed: two concepts with same wid)
  \item NEs not tagged (try to do at least some)
  \item Easy online lookup
    \\ by word, lemma, POS, synset, hypernym, file
  \item New words (from tags 'u', 's', new wordnets, pron, int, dates)
    \begin{itemize}
    \item Pronouns, question words not tagged:   add by hand
    \item Need tool for new entries
    \end{itemize}
  \end{itemize}
\item Crosslingual annotation (Summer)
  \begin{itemize}
  \item Goal (14,000 sentence pairs)
    \begin{itemize}
    \item C-E 6000 (1200 done)
    \item M-E 2000
    \item J-E 6000
    \end{itemize}
  \item Note: News corpus aligned by NICT: hire student to convert? (4,000!)
  \end{itemize}
\item Papers/Research
  \begin{itemize}
  \item  Comparison of WN cover, estimate of missing words/senses
    \\ in the style of \citet{Baldwin:Bender:Flickinger:Kim:Oepen:2004}
    \\ for the four genre: how much missing, how much overlap
    \\ /home/bond/papers/wn/wn-ml-tag
    \\ ACL short, LREC?
  \item xling annotation paper (based on Gao (FYP) and \citet{Mok:2012})
    \\ PACLIC?, ACL short?
\item Translation Shifts (main research: ACL next year, CL)
\item Transfer Rule Acquisition
\item Xling WSD/Annotation
\item FYP: Pronouns
\item Idioms
  \end{itemize}
\item Misc
  \begin{itemize}
  \item Chinese Wordnet --- release?
    \begin{itemize}
    \item Chinese WN
    \item Gloss corpus
    \item Classifiers
    \item[?]
    \end{itemize}
  \item ZhEn 
  \end{itemize}

\end{itemize}


\section{Guidelines for extending the wordnets}

Our goal in building the wordnet is to have coverage over all open
class predicates, whether from single words, multi-word expressions or
constructions.  We do not attempt to cover all proper names (entities)
although we have some common ones.  We also do not cover closed class
mainly functional words.  We differ from the Princeton wordnet in
adding pronouns, as we wish to be able to compare the use of pronouns
with the use of other referring expressions.

\subsection{What goes in}

\subsubsection{Lemmas (Lexical Units)}


\subsubsection{Pronouns}

\subsubsection{Interrogatives}

\subsection{What stays out}





\chapter{Applications}

What we do with the corpus.

\section{Wordnet Frequency and Examples}

\section{On-line look-up}



\subsection*{Acknowledgments}

Much of the initial work on adding senses was done by Jeanette Tan,
Eshley Gao, Hazel Mok and Yue Hui Ting, supported by the NTU URECA
program.  We also received support from NICT, in the form of
internships and Toyohashi University of Technology, in the person of
Takayuki Kuribayashi.  Some annotation was carried out at NTT-AT under
the capable supervision of Satsuki Abe.

The corpus has been used to teach lexical semantics at NTU in both
HG2002 \textit{Semantics and Pragmatics} and HG250 \textit{Chinese
  Lexicology}.  Thanks to the students for their comments and
feedback.  

Mathieu Morey was deeply involved in discussions of how to represent
the cross-lingual alignment.

We received much useful advice from members of the KYOTO project,
especially Christiane Fellbaum, Piek Vossen and Hitoshi Isahara. 

Kyoko Kanzaki and Hitoshi Isahara helped extend the Japanese
translation of the Cathedral and the Bazaar.  Kyonghee Paik helped
sentence aligning the Speckled Band.

Funding: JSPS, SUG, Tier 1, NICT.

\bibliographystyle{plainnat}
\bibliography{abb,mtg,nlp,ling}


\appendix
\chapter{Part of Speech Tags}
\section{English Part of Speech Tags}

We use the Penn POS tags with the addition of \tg{VAX}

\begin{table}
  \centering
  \begin{tabular}{lll}
    \textbf{Tag} & \textbf{Description} & \textbf{Example} & \textbf{Comment}\\
\hline
%1.
\tg{CC} & Coordinating conjunction \\
%2.
\tg{CD} & Cardinal number \\
%3.
\tg{DT} & Determiner \\
%4.
\tg{EX} & Existential there \\
%5.
\tg{FW} & Foreign word \\
%6.
\tg{IN} & Preposition or subordinating conjunction \\
%7.
\tg{JJ} & Adjective \\
%8.
\tg{JJR} & Adjective, comparative \\
%9.
\tg{JJS} & Adjective, superlative \\
%10
\tg{LS} & List item marker \\
%11
\tg{MD} & Modal \\
%12
\tg{NN} & Noun, singular or mass \\
%13
\tg{NNS} & Noun, plural \\
%14
\tg{NNP} & Proper noun, singular \\
%15
\tg{NNPS} & Proper noun, plural \\
%16
\tg{PDT} & Predeterminer \\
%17
\tg{POS} & Possessive ending \\
%18
\tg{PRP} & Personal pronoun \\
%19
\tg{PRP\$} & Possessive pronoun \\
%20
\tg{RB} & Adverb \\
%21
\tg{RBR} & Adverb, comparative \\
%22
\tg{RBS} & Adverb, superlative \\
%23
\tg{RP} & Particle \\
%24
\tg{SYM} & Symbol \\
%25
\tg{TO} & to \\
%26
\tg{UH} & Interjection \\
%27
\tg{VB} & Verb, base form \\
%28
\tg{VBD} & Verb, past tense \\
%29
\tg{VBG} & Verb, gerund or present participle \\
%30
\tg{VBN} & Verb, past participle \\
%31
\tg{VBP} & Verb, non-3rd person singular present \\
%32
\tg{VBZ} & Verb, 3rd person singular present \\
% New
\rowcolor{lightgray} 
\tg{VAX} & Verb, Auxiliary & \eng{have, be} & Don't tag\\
%33
\tg{WDT} & Wh-determiner \\
%34
\tg{WP} & Wh-pronoun \\
%35
\tg{WP\$} & Possessive wh-pronoun \\
%36
\tg{WRB} & Wh-adverb  \\

  \end{tabular}
  \caption{English POS Tags}
  \label{tab:eng-pos}
\end{table}

\section{Japanese Part of Speech Tags}

We use the IPA POS tags with the addition of \tg{動詞-変接続}
%%% Taken from http://www.unixuser.org/~euske/doc/postag/
\begin{verbatim}
ChaSen 品詞体系 (IPA品詞体系)

ChaSen の品詞体系は任意の階層化を許している。 いわゆる形容動詞は名詞の形容動詞語幹として含まれ、 形容詞には含まれない。Juman の指示詞という カテゴリは「連体詞」に含まれている。 判定詞「だ」は助動詞とされている。

Type1 	Type2 	Type3 	Type4 	Examples 	Description
連体詞 	「この」、「いろんな」、「おっきな」、「堂々たる」 	
接頭詞 	形容詞接続 	「お」、「まっ」、「クソ」 	「お高い」「クソ高い」「まっ赤」
数接続 	「No.」、「計」、「毎分」 	あとに数値がくるもの。
動詞接続 	「ぶっ」、「引き」 	あとに動詞がくるもの。「ぶったたく」
名詞接続 	「もと」、「アンチ」、「最」、「総」 	あとに名詞がくるもの。
名詞 	引用文字列 	「いわく」のみ 	
サ変接続 	「苦労」、「終了」、「アピール」、「くしゃみ」 	「～する」の形をとれるもの。
ナイ形容詞語幹 	「申し訳」、「とんでも」、「おとなげ」 	「～ない」の形をとれるもの。
形容動詞語幹 	「あからさま」、「ミステリアス」、「決定的」、「無人」 	「～な」の形をとれるもの。
動詞非自立的 	「ごらん」、「ちょうだい」のみ 	「～してごらん」「～してちょうだい」
副詞可能 	「１０月」、「せんだって」、「当分」 	おもに時相名詞。
一般 	「大根」、「シエスタ」、「加速度」、「ありさま」 	一般名詞。
数 	「ゼロ」、「億」 	数値。
接続詞的 	「VS」、「対」、「兼」のみ 	
固有名詞 	一般 	「北穂高岳」、「電通銀座ビル」、「Ｇ１」 	下のカテゴリにあてはまらないもの。山・川・商品名など。
人名 	一般 	「グッチ裕三」、「紫式部」 	芸能人名など。
姓 	「山田」、「ビスコンティ」 	
名 	「Ｂ作」、「アントニオ」、「右京太夫」 	
組織 	「いすゞ自動車」、「ニチレイ」、「統一アイルランド党」 	
地域 	一般 	「北海道」、「やながわ工業団地」、「ラムサール」 	
国 	「露西亜」、「バングラデシュ」 	
接尾 	サ変接続 	「化」、「入り」 	「巨大化」「キャンプ入り」
一般 	「ぎみ」、「ゆかり」、「枚」、「不可」 	
形容動詞語幹 	「がち」、「的」、「同然」 	
助数詞 	「オクターブ」、「％」、「ヶ国」 	単位。
助動詞語幹 	「そ」、「そう」のみ 	「～してそう」「～そうだ」で使われる。
人名 	「君」、「はん」 	
地域 	「州」、「市内」、「港」 	
特殊 	「ぶり」、「み」、「方」 	
副詞可能 	「いっぱい」、「前後」、「次第」 	
代名詞 	一般 	「そこ」、「俺」、「こんちくしょう」 	
縮約 	「わたしゃ」、「そりゃあ」 	
非自立 	一般 	「こと」、「きらい」、「くせ」、「もの」 	
形容動詞語幹 	「ふう」、「みたい」のみ 	
助動詞語幹 	「よ」、「よう」のみ 	
副詞可能 	「限り」、「さなか」、「うち」 	
特殊 	助動詞語幹 	「そ」、「そう」のみ 	
動詞 	自立 	「いがみ合う」、「たてつく」、「垢抜ける」 	
接尾 	「する」、「られる」、「させる」、「がかる」 	
非自立 	「しまう」、「ちゃう」、「願う」 	「行ってしまう」「やっちゃったね」「ご遠慮願う」
形容詞 	自立 	「けたたましい」、「分別臭い」、「めでたい」 	
接尾 	「ったらしい」、「っぽい」 	
非自立 	「づらい」、「がたい」、「よい」 	「聞きづらい」「見よい」
副詞 	一般 	「たいそう」、「人一倍」、「いけしゃあしゃあ」 	
助詞類接続 	「あまり」、「いつも」、「ぱさぱさ」 	同一文節内に付属語がこれるもの。
接続詞 	「けれども」、「たとえば」、「反面」 	
助詞 	格助詞 	一般 	「の」、「から」、「を」 	
引用 	「と」のみ 	
連語 	「について」、「とかいう」 	解析ミスを防ぐためのマクロ。
係助詞 	「は」、「こそ」、「も」、「や」 	
終助詞 	「かしら」、「ぞ」、「っけ」、「わい」 	
接続助詞 	「て」、「つつ」、「および」、「ので」 	
特殊 	「かな」、「けむ」、「にゃ」 	「せにゃならん」 これ以外の用途は不明。
副詞化 	「と」、「に」のみ 	
副助詞 	「くらい」、「なんか」、「ばっかり」 	
副助詞／並立助詞／終助詞 	「か」のみ 	
並立助詞 	「とか」、「だの」、「やら」 	
連体化 	「の」のみ 	
助動詞 	「らしい」、「ござる」、「っす」、「じゃん」 	「そうでございます」
感動詞 	「うむ」、「お疲れさま」、「トホホ」 	
記号 	句点 	「。」、「．」のみ 	
読点 	「、」、「，」のみ 	
空白 	「　」のみ 	
アルファベット 	「Ａ」 	
一般 	「？」、「！」、「￥」 	
括弧開 	「（」、「【」 	
括弧閉 	「”」、「＞」 	
フィラー 	「えーと」、「なんか」 	(たぶん)音声認識用。
その他 	間投 	「あ」、「ア」のみ 	「そんなぁ」
未知語 	「オマエモナー」「ＮＵＬＬＰＯ」 	
\end{verbatim}




\end{CJK}
\end{document}


Notes: try to revise as we go
 * 

 * recommend chrome/firefox

 * can do batch updates
 * try to link to wn-feedback

 * harder to access next week (but possible)
 * try to link as much as possible
 * linking Pronouns

 * report once a week (Thursday)

 